\section{Stakeholders e necessidades}
\begin{table}[!ht]\centering\scriptsize\renewcommand{\arraystretch}{1.35}
\begin{tabularx}{\linewidth}{|>{\bfseries\raggedright\arraybackslash}p{2.5cm}|>{\raggedright\arraybackslash}X|>{\raggedright\arraybackslash}X|>{\raggedright\arraybackslash}X|}
\hline
\textbf{Stakeholder / papel} & \textbf{Participação} & \textbf{Necessidade identificada} & \textbf{Origem / evidência} \\
\hline
Emissor do aviso &
Elabora comunicados, seleciona público e canais, solicita validação/distribuição e publica &
Criar comunicados estruturados, efetivos e seguros em relação à autorização prévia &
BPMN: criar, selecionar, solicitar autorização, enviar/postar e cancelar avisos. \\
\hline
Avaliador &
Analisa solicitações, autoriza, solicita ajustes ou indefere avisos &
Moderar o conteúdo divulgado, bloquear informações inviáveis e direcionar adequadamente &
BPMN: avaliar solicitação; comunicar aprovação ou inviabilidade motivada. \\
\hline
Destinatários &
Recebem notificações e consultam informações &
Acessar avisos relevantes no tempo adequado, sem enxurrada de comunicados aleatórios &
BPMN: receber o aviso, comunicado ou resposta. \\
\hline
Tutor do aluno &
Inicia o fluxo enviando uma mensagem particular e recebe a resposta final &
Estabelecer canal oficial e documentado para resolver demandas &
BPMN: criar, enviar, receber resposta da mensagem. \\
\hline
Respondente &
Avalia mensagens, redistribui ou responde diretamente &
Filtrar, avaliar e dar retorno adequado às demandas de forma ágil &
BPMN: avaliar, devolver e responder mensagens. \\
\hline
\end{tabularx}
\end{table}

As necessidades são evidenciadas pelo BPMN e pelo processo AS-IS: o emissor precisa estruturar o conteúdo; o avaliador precisa controlar a publicação; e os destinatários precisam receber apenas os comunicados pertinentes ao seu contexto.
